% ======================================================================
% sec-consequences.tex -- consequences of the filtration:
% quantitative comparison, commitment responsibility, a bounded-treewidth
% algorithm (in the body), and the PDB connection.
% ======================================================================
\section{Consequences of the Filtration}
\label{sec:discussion}

When a domain's conflict or dependency structure fixes a natural commitment
basis, different specifications over that basis share a single filtration:
their resolutions are selected from one ambient space $\Dets$, and support
depth is compared over it. For $\mbox{Datalog}^\neg$ the negation semantics
are different reading depths of this shared filtration
(Theorem~\ref{thm:negation-general}); for transactions the isolation levels
are different resolving subsets of it (Appendix~\ref{sec:transactions}). We
draw three consequences: quantitative comparison across specifications,
responsibility attribution, and an algorithm.
Appendix~\ref{app:worked-queries} runs the whole query family---support, depth,
composition, cross-semantics comparison, responsibility, and the re-enabled
$K$-provenance---end to end on the running example.

\subsection{Work Regret and Semantic Shift}
\label{sec:regret-shift}

Fix a provenance query and an observable $a$ with support depth
$\sdepth_I(a)$ under specification $I$. Comparing a workload under $I$ against
a candidate alternative $I'$ over the same basis, the difference decomposes.

\paragraph{Work regret.}
If $a$ is robust under both $I$ and $I'$ but $\sdepth_I(a) > \sdepth_{I'}(a)$,
the extra commitment layers under $I$ did not change the answer---\emph{work
regret}, resolution spent to no semantic effect. In $\mbox{Datalog}^\neg$,
stable-model evaluation of an atom that is already settled by the stratified
prefix regrets its choice layers (and the transaction instantiation of
Appendix~\ref{sec:transactions} exhibits regret in both directions between
isolation levels).

\paragraph{Semantic shift.}
If $a$'s why-support itself differs between $I$ and $I'$---robust under one,
contingent under the other---the meaning genuinely changes. Moving from
stable-model to well-founded semantics converts choice-layer atoms to
$\mathbf{u}$ (Theorem~\ref{thm:negation-general}); the transaction
instantiation witnesses the same phenomenon across isolation levels
(Appendix~\ref{sec:transactions}).

\begin{proposition}[Compositional depth bound]
  \label{prop:depth-bound}
  If every input of a positive relational provenance query has support depth
  $< k$, so does every output
  (Theorem~\ref{thm:filtration-respected}): layers $k, k{+}1, \dots$ are inert
  for that query.
\end{proposition}

\noindent
Certifying this exactly is coNP-complete in the number of commitments in the
relevant layers (Appendix~\ref{app:robustness-proofs}), but the lead
instantiation admits a polynomial sufficient check: in $\mbox{Datalog}^\neg$,
absence of transitive dependency on any negative SCC~\cite{gebser2012asp}
guarantees support depth below the choice layers. (The transaction
instantiation admits an analogous check; Appendix~\ref{sec:transactions}.)

\subsection{Commitment Responsibility}
\label{sec:responsibility}

Work regret finds wasted layers; \emph{responsibility} finds the single
commitment most accountable for an observable's contingency. For a commitment
occurrence $e$ recorded in the traces, let $\Dets_e \subseteq \Dets$ be the
determinations that commit $e$, and let $\operatorname{Piv}_a(e) \subseteq
\Dets_e$ be those whose support answer for $a$ flips when $e$ is withheld.

\begin{definition}[Responsibility value]
  \label{def:responsibility}
  The \emph{responsibility} of $e$ for $a$ is the density of its pivotal set
  among the determinations that commit it,
  \[
    \operatorname{resp}_a(e)
    = \frac{|\operatorname{Piv}_a(e)|}{|\Dets_e|}
    \qquad (\Dets_e \neq \emptyset),
  \]
  under the counting measure on a finite $\Dets_e$ (or the conditional density
  $\mu(\operatorname{Piv}_a(e) \mid \Dets_e)$ under a measure $\mu$).
\end{definition}

\noindent
Responsibility $0$ means $e$ is never pivotal; $1$ means every determination
committing $e$ has its answer settled by $e$. This is a determination-level
query of the kind the filtration already handles: its value is fixed by any
canonical prefix that fixes the pivotal set, so its certificate depth equals
that of $\mathbf{1}_{\operatorname{Piv}_a(e)}$. It is natural to refine this
pivotal-set density into a \emph{Shapley value} over commitment
occurrences---the sequential counterpart of tuple-level
responsibility~\cite{meliou2010causality,livshits2021shapley}, with semantic
commitments as players in place of data tuples. We leave that development,
including its hardness and bounded-treewidth tractability, to the
self-contained Appendix~\ref{app:responsibility-details}.

\subsection{An Algorithm: Support and Robustness at Bounded Treewidth}
\label{sec:algorithm}

Enumerating $\Dets$ is exponential in general, so support and robustness need
a structural handle. Within a single commitment layer the commitments commute,
and a support is a positive Boolean formula over commitment variables (the
depth-$1$ case of Proposition~\ref{prop:pdb}). Robustness---is the support
all of $\Dets$?---is then validity of that formula, and support membership is
model counting. Both are tractable when the formula has bounded treewidth.

\begin{theorem}[Bounded-treewidth support and robustness]
  \label{thm:treewidth}
  Fix a provenance query $Q$ and an output observable $a$ whose single-layer
  support formula $\phi_a$ (over $n$ commitment variables) has primal-graph
  treewidth $w$. Then:
  \begin{enumerate}[nosep,label=(\alph*)]
    \item the why-support of $a$, and hence robustness
          ($Y_a = \Dets$) and impossibility ($Y_a = \emptyset$), can be
          decided in time $O(2^w \cdot \mathrm{poly}(n))$;
    \item under a product measure on commitment variables, the support
          density $\mu(Y_a)$ is computable in the same time by weighted model
          counting on a tree decomposition.
  \end{enumerate}
\end{theorem}
\begin{proof}[Proof sketch]
  By Proposition~\ref{prop:support-hom} the support of a positive-RA query is
  obtained from base supports by $\cup$ and $\cap$, so $\phi_a$ is a monotone
  Boolean combination of the base commitment variables; its primal graph has
  treewidth $w$ by hypothesis. Validity, satisfiability, and (weighted) model
  counting on a bounded-treewidth formula run in $O(2^w \cdot \mathrm{poly}(n))$
  by dynamic programming over a tree decomposition. Robustness is validity of
  $\phi_a$; impossibility is unsatisfiability; density is the weighted model
  count. Full statement and proof in Appendix~\ref{app:robustness-proofs}.
\end{proof}

\noindent
The treewidth here is that of the \emph{support formula}, not of the
underlying conflict or dependency graph; the two coincide under a locality
condition (each clause involves only commitments sharing a dependency edge)
that both instantiations' natural workloads satisfy
(Appendix~\ref{app:robustness-proofs}). Multi-layer supports are handled layer
by layer: Theorem~\ref{thm:filtration-respected} bounds the depth that must be
examined, and within each layer the single-layer algorithm applies.

\subsection{Connection to Probabilistic Databases}
\label{sec:pdb}

The filtration locates determination provenance precisely within the
probabilistic-database complexity landscape.

\begin{proposition}[Depth-$1$ is tuple-independent]
  \label{prop:pdb}
  A depth-$1$ determination semiring over $n$ binary commitments is isomorphic
  to a tuple-independent PDB over $n$ existence events: supports are positive
  Boolean formulas, and query evaluation has the same complexity.
\end{proposition}

\noindent
At higher depths the filtration structures the correlations: within a layer
commitments commute (PDB-style lineage), while cross-layer dependencies
introduce conditioning. Determination provenance thus sits between
tuple-independent PDBs (depth $1$) and general correlated PDBs (intractable):
the PDB dichotomy~\cite{dalvi2012probabilistic} applies \emph{within} each
layer, and support depth measures how many layers must be composed. Under a
probability measure on $\Dets$, support densities become probabilities and
Theorem~\ref{thm:treewidth} becomes a weighted-model-counting bound of exactly
the PDB kind.
